\section{Preliminaries}
\label{sec:preliminaries}

This section fixes the objects used throughout the paper. We separate three levels which are often conflated in branching arguments: the deterministic Duhamel expansion, the intrinsic signed measure attached to a marked genealogy, and a positive proposal used to sample that signed measure. The distinction is essential for both the positive and negative results.

\subsection{Heat semigroup and mild equations}

Write
\begin{equation*}
  \Torus=\R/(2\pi\mathbb Z)
\end{equation*}
for the one-dimensional torus. Let $(P_t)_{t\geq0}$ be the periodic heat semigroup generated by $\frac12\partial_x^2$. Thus, for a bounded measurable $f$,
\begin{equation*}
  P_tf(x)=\int_{\Torus}p_t^{\Torus}(x,y)f(y)\,dy,
\end{equation*}
where $p_t^{\Torus}$ is the periodic heat kernel. We also use the same notation for the heat semigroup on $\R$ when the spatial domain is clear from context.

The general terminal equation considered in the coding-tree part is
\begin{equation}
  \partial_tu(t,x)
  +\frac12\partial_x^2u(t,x)
  +f(J_nu(t,x))=0,
  \qquad
  u(T,x)=\phi(x),
  \label{eq:general-terminal-pde}
\end{equation}
with finite spatial jet
\begin{equation*}
  J_nu=(u,u_x,\ldots,u_{x^n}).
\end{equation*}
For smooth solutions, Duhamel's formula reads
\begin{equation*}
  u(t)
  =
  P_{T-t}\phi
  +\int_t^T
  P_{s-t}\br{f(J_nu(s))}\,ds.
\end{equation*}
The coding-tree construction of Nguwi, Penent, and Privault~\cite{NguwiPenentPrivault2023} is built from repeated differentiation and substitution in this identity.

The main equation of the paper is the forward quadratic Hessian problem
\begin{equation}
  \partial_tv
  =
  \frac12v_{xx}+\lambda(v_{xx})^2,
  \qquad
  v(0)=\phi,
  \label{eq:hessian-forward}
\end{equation}
on $[0,T]\times\Torus$. Put
\begin{equation*}
  z=v_{xx},
  \qquad
  L(t)=P_t\phi''.
\end{equation*}
Then
\begin{equation}
  z(t)
  =
  L(t)
  +\lambda\int_0^t
  \partial_x^2P_{t-s}\br{z(s)^2}\,ds.
  \label{eq:hessian-mild}
\end{equation}
Define the Hessian Duhamel operator
\begin{equation}
  (\cD f)(t)
  =
  \int_0^t\partial_x^2P_{t-s}f(s)\,ds.
  \label{eq:def-D}
\end{equation}
Then~\eqref{eq:hessian-mild} is $z=L+\lambda\cD(z^2)$.

\subsection{Centered Gaussian derivative marks}

Let $Z\sim N(0,1)$ and let
\begin{equation*}
  \He_2(z)=z^2-1
\end{equation*}
be the second probabilists' Hermite polynomial. Gaussian integration by parts gives, for smooth $F$ and $r>0$,
\begin{equation}
  \partial_x^2P_rF(x)
  =
  \frac1r\E\br{\He_2(Z)F(x+\sqrt rZ)}.
  \label{eq:gaussian-hessian-unsubtracted}
\end{equation}
Since $\E\He_2(Z)=0$, the same identity may be written in centered form,
\begin{align}
  \partial_x^2P_rF(x)
  &=
  \E\br{\widehat K_rF(x,Z)},
  \label{eq:centered-hessian-mark}\\
  \widehat K_rF(x,Z)
  &=
  \frac{\He_2(Z)}r
  \br{F(x+\sqrt rZ)-F(x)}.
  \label{eq:centered-hessian-random-field}
\end{align}
The subtraction does not change the signed expectation. It is essential in absolute estimates. If $0<\alpha<1$, then
\begin{equation}
  \E\abs{\widehat K_rF(x,Z)}
  \leq
  c_{2,\alpha}r^{-1+\alpha/2}[F]_{C^\alpha},
  \qquad
  c_{2,\alpha}
  =
  \E\br{\abs{\He_2(Z)}\abs Z^\alpha}.
  \label{eq:centered-hessian-holder}
\end{equation}
The short-edge singularity in~\eqref{eq:centered-hessian-holder} is integrable in $r$. By contrast, if one propagates the pathwise $C^\alpha$ seminorm of the random field $x\mapsto\widehat K_rF(x,Z)$, the direct estimate has scale $r^{-1}$ and no H\"older gain. This observation motivates the later decision to average several signed marks before taking the first absolute value.

For a chain of independent Gaussian increments $Z_1,\ldots,Z_m$ with durations $r_1,\ldots,r_m$, put
\begin{equation*}
  Y=\sum_{j=1}^m\sqrt{r_j}Z_j,
  \qquad
  R=\sum_{j=1}^m r_j.
\end{equation*}
Conditioning the product of second Hermite scores on $Y$ gives the bridge-score identity
\begin{equation}
  \E\br{
    \left.
    \prod_{j=1}^m\frac{\He_2(Z_j)}{r_j}
    \right|Y
  }
  =
  \frac{\He_{2m}(Y/\sqrt R)}{R^m}.
  \label{eq:bridge-score}
\end{equation}
For a bare derivative chain,~\eqref{eq:bridge-score} is the basic example of cancellation obtained by averaging internal Gaussian bridge coordinates before an absolute value. Standard Hermite identities may be found, for example, in~\cite{Janson1997}.

\subsection{Signed measures and positive proposals}

The measure-theoretic formulation makes proposal invariance transparent. Let $(M,\mathcal A)$ be a measurable mark space and let $\mu$ be a finite signed measure on $M$. Its total-variation norm is denoted by $\norm\mu_{\TV}$.

If $Q$ is a probability measure such that $\mu\ll Q$, the canonical importance-sampling weight is
\begin{equation*}
  W=\frac{d\mu}{dQ}.
\end{equation*}
Then
\begin{equation}
  \int_M W\,dQ
  =
  \mu(M),
  \qquad
  \int_M\abs W\,dQ
  =
  \norm\mu_{\TV}.
  \label{eq:proposal-tv}
\end{equation}
Thus a proposal changes the density of the estimator but cannot change the $L^1$ cost of sampling a fixed signed measure exactly. The same conclusion holds if the proposal variables are dependent: $Q$ is simply their joint law.

Auxiliary unbiased randomness cannot reduce the first moment if the canonical signed density remains its conditional barycenter. More precisely, let $R$ denote the exposed canonical mark coordinate and let $Y\in L^1$ satisfy
\begin{equation}
  \E_Q[Y\mid\sigma(R)]
  =
  \frac{d\mu}{dQ}(R).
  \label{eq:conditional-barycenter-general}
\end{equation}
Conditional Jensen gives
\begin{equation}
  \E_Q\abs Y
  \geq
  \E_Q\abs{\E_Q[Y\mid\sigma(R)]}
  =
  \norm\mu_{\TV}.
  \label{eq:conditional-jensen-tv}
\end{equation}
We always use~\eqref{eq:conditional-barycenter-general} under the assumption that $Y\in L^1$, typically in a proof by contradiction. We never define an unresolved nonintegrable object by an ordinary conditional expectation.

The capstone theorem also uses conditional expectation with respect to a finite positive measure $\nu$ which need not be normalized. For $G\in L^1(\nu)$ and a sub-sigma-field $\mathcal G$, we write
\begin{equation*}
  \E_\nu[G\mid\mathcal G]
\end{equation*}
for the $\mathcal G$-measurable function characterized by
\begin{equation*}
  \int_A\E_\nu[G\mid\mathcal G]\,d\nu
  =
  \int_A G\,d\nu,
  \qquad A\in\mathcal G.
\end{equation*}
When $\nu(M)>0$, this is the usual conditional expectation after normalizing $\nu$ to a probability measure; the normalization cancels from the defining identity.

\begin{definition}[Raw-faithfulness]
\label{def:raw-faithful}
Let $\mu$ be the intrinsic signed measure on a canonical raw marked state space, and let $Q$ be a probability proposal dominating it. An $L^1(Q)$ estimator is \emph{raw-faithful} if its conditional barycenter after the canonical raw state is exposed is $d\mu/dQ$, as in~\eqref{eq:conditional-barycenter-general}.
\end{definition}

The definition permits arbitrary changes of positive proposal and arbitrary auxiliary conditionally unbiased randomization. It excludes transformations which move signed mass between different canonical raw states before the conditional barycenter is taken.

\subsection{Decorative and non-barycentric randomness}

Raw-faithfulness is stronger than the statement that a simulation contains random marks. The distinction is elementary but important.

Suppose $S$ is a discrete skeleton and $F_S$ is a deterministic value. Sampling an independent Gaussian family $G$ and returning
\begin{equation*}
  Y=\frac{F_S}{\pi(S)}
\end{equation*}
produces a simulation state which contains continuous randomness, but the Gaussian variables are decorative. They can be deleted without changing the estimator. A theorem which only prohibits ``random marks'' would therefore be ill posed.

At the other extreme, a random mark can matter while raw-faithfulness fails. The antithetic second-derivative estimator
\begin{equation}
  \widetilde K_rF(x,Z)
  =
  \frac{\He_2(Z)}{2r}
  \br{F(x+\sqrt rZ)+F(x-\sqrt rZ)-2F(x)}
  \label{eq:antithetic-prelim}
\end{equation}
has the same signed expectation as~\eqref{eq:centered-hessian-mark}, but its conditional value at a fixed $Z$ is different. It redistributes signed mass between the raw states $Z$ and $-Z$. The negative theorem in~\cref{sec:raw-faithful} does not apply to~\eqref{eq:antithetic-prelim}.

\subsection{Binary Duhamel trees and patches}

Let $\cT$ denote the set of finite rooted planar full binary trees. A leaf represents the heat lift $L=P_\cdot\phi''$. If $\tau=[\tau_1,\tau_2]$ has left and right subtrees $\tau_1$ and $\tau_2$, its deterministic Duhamel contribution is defined recursively by
\begin{equation}
  F_\bullet=L,
  \qquad
  F_{[\tau_1,\tau_2]}
  =
  \lambda\cD(F_{\tau_1}F_{\tau_2}).
  \label{eq:tree-profile-recursion}
\end{equation}
For the finite Picard expansion, these are finite iterated Duhamel integrals.

A \emph{patch} is a maximal chain of internal vertices obtained by repeatedly following the left child. Every internal vertex belongs to exactly one patch. Contracting the maximal left chains gives the \emph{patch skeleton}. A \emph{decorated patch skeleton} records both the length of every contracted chain and the ordered side-subtree attachment slots along that chain. With this decoration, finite planar full binary trees and decorated patch skeletons are in bijection; the precise inverse is proved in~\cref{prop:tree-patch-bijection} below.

For a patch of length $m$, with side profiles $b_1,\ldots,b_m$ and leftmost terminal profile $G$, its complete signed contribution is
\begin{equation}
  \cP_m[b_1,\ldots,b_m;G](t)
  =
  \lambda^m
  \int_{0<s_1<\cdots<s_m<t}
  \partial_x^2P_{t-s_m}\Xi_m
  \,ds_1\cdots ds_m,
  \label{eq:complete-patch}
\end{equation}
where
\begin{align*}
  \Xi_1&=b_1(s_1)P_{s_1}G,\\
  \Xi_r&=b_r(s_r)\partial_x^2P_{s_r-s_{r-1}}\Xi_{r-1},
  \qquad 2\leq r\leq m.
\end{align*}
We set $\cP_0[G](t)=P_tG$.

\subsection{Canonical raw signed measures for finite Hessian trees}
\label{subsec:canonical-raw-measures}

The later notions of raw-faithfulness, coarsening, and residual signed variation require a genuine finite signed measure for each fixed finite Duhamel tree. We construct that measure here. The construction is finite-depth and uses only one positive spatial H\"older exponent at the terminal Hessian datum.

Fix
\begin{equation*}
  0<\alpha<1,
  \qquad
  T>0,
  \qquad
  g=\phi''\in C^\alpha(\Torus).
\end{equation*}
No smallness assumption on $g$ or $\lambda$ is imposed in this subsection. Write $\gamma$ for standard Gaussian measure on $\R$.

\paragraph{Raw mark spaces.}
For the leaf tree $\bullet$, set
\begin{equation*}
  \Omega_\bullet=\R,
  \qquad
  \nu_\bullet=\gamma.
\end{equation*}
If $\tau=[\tau_1,\tau_2]$, define recursively
\begin{equation}
  \Omega_\tau
  =
  [0,T]\times\R\times
  \Omega_{\tau_1}\times\Omega_{\tau_2},
  \qquad
  \nu_\tau
  =
  ds\otimes\gamma(dz)\otimes
  \nu_{\tau_1}\otimes\nu_{\tau_2}.
  \label{eq:raw-space-recursion}
\end{equation}
We equip every $\Omega_\tau$ with its product Borel sigma-field. Thus every $\Omega_\tau$ is standard Borel and every $\nu_\tau$ is finite. If $|\tau|$ is the number of internal vertices, then
\begin{equation*}
  \nu_\tau(\Omega_\tau)=T^{|\tau|}.
\end{equation*}
The chronological constraints are imposed by the density below rather than by changing the measurable space with the observation time.

\paragraph{Canonical centered raw density.}
For a leaf, define
\begin{equation}
  \mathscr R_\bullet^{t,x}(\xi)
  =
  g(x+\sqrt t\,\xi).
  \label{eq:raw-density-leaf}
\end{equation}
Spatial addition is understood modulo $2\pi$. If $\tau=[\tau_1,\tau_2]$ and
\begin{equation*}
  \omega=(s,z,\omega_1,\omega_2)\in\Omega_\tau,
\end{equation*}
then for $0<s<t$ put
\begin{equation*}
  r=t-s,
  \qquad
  x_z=x+\sqrt r\,z,
\end{equation*}
and set
\begin{align}
  \mathscr R_\tau^{t,x}(\omega)
  ={}&
  \lambda\frac{\He_2(z)}r
  \Bigl[
    \mathscr R_{\tau_1}^{s,x_z}(\omega_1)
    \mathscr R_{\tau_2}^{s,x_z}(\omega_2)
  \notag\\
  &\hspace{7em}-
    \mathscr R_{\tau_1}^{s,x}(\omega_1)
    \mathscr R_{\tau_2}^{s,x}(\omega_2)
  \Bigr].
  \label{eq:raw-density-internal}
\end{align}
For $s\geq t$, including $s=t$, define $\mathscr R_\tau^{t,x}=0$. The shifted and unshifted terms in~\eqref{eq:raw-density-internal} use the same descendant marks $(\omega_1,\omega_2)$. Thus~\eqref{eq:raw-density-internal} is exactly the centered one-edge realization~\eqref{eq:centered-hessian-random-field} applied to the product of two independent descendant raw fields.

By induction on $\tau$, the map
\begin{equation*}
  (t,x,\omega)
  \longmapsto
  \mathscr R_\tau^{t,x}(\omega)
\end{equation*}
is jointly Borel measurable on $[0,T]\times\Torus\times\Omega_\tau$. The leaf map is continuous. At an internal vertex, the region $\{0<s<t\}$ is Borel, the maps $r=t-s$ and $x+\sqrt r z$ are Borel there, and~\eqref{eq:raw-density-internal} is obtained from the child densities by measurable algebraic operations. Defining the value to be zero on $s\geq t$ removes the diagonal singularity from the measurable definition.

\paragraph{$L^1$-valued H\"older norm.}
Let $(E,\mathcal E,m)$ be a finite positive measure space. For a jointly measurable field $H:\Torus\times E\to\R$ and $0<\eta<1$, define
\begin{align}
  \norm{H}_{\mathcal H^\eta(m)}
  ={}&
  \sup_{x\in\Torus}\int_E\abs{H(x,e)}\,dm(e)
  \notag\\
  &+
  \sup_{x\neq y}
  \frac{
    \displaystyle
    \int_E\abs{H(x,e)-H(y,e)}\,dm(e)
  }{d_\Torus(x,y)^\eta}.
  \label{eq:L1-valued-holder}
\end{align}
This is the $C^\eta$ norm of the map $x\mapsto H(x,\cdot)$ with values in $L^1(m)$. In particular, it does not take a pathwise H\"older norm before averaging.

For $r>0$ define
\begin{equation*}
  (\mathcal K_rH)(x,z,e)
  =
  \frac{\He_2(z)}r
  \br{H(x+\sqrt r z,e)-H(x,e)}.
\end{equation*}

\begin{lemma}[One centered edge with strict regularity loss]
\label{lem:raw-one-edge-loss}
Let $0<\beta<\eta<1$ and $T>0$. There is $C_{\beta,\eta,T}<\infty$ such that, for $0<r\leq T$,
\begin{equation}
  \norm{\mathcal K_rH}_{\mathcal H^\beta(\gamma\otimes m)}
  \leq
  C_{\beta,\eta,T}
  r^{-1+(\eta-\beta)/2}
  \norm H_{\mathcal H^\eta(m)}.
  \label{eq:raw-one-edge-loss}
\end{equation}
Consequently the right side is integrable in $r$ on $(0,T)$.
\end{lemma}

\begin{proof}
For the supremum part, the input H\"older increment gives
\begin{equation*}
  \sup_x\int\abs{\mathcal K_rH(x,z,e)}
  \,d\gamma(z)dm(e)
  \leq
  c_{2,\eta}r^{-1+\eta/2}
  \norm H_{\mathcal H^\eta(m)}.
\end{equation*}
For the spatial increment, write $d=d_\Torus(x,y)$. Applying the input H\"older estimate to both translated terms gives
\begin{equation}
  \int
  \abs{\mathcal K_rH(x)-\mathcal K_rH(y)}
  \leq
  C r^{-1}d^\eta\norm H_{\mathcal H^\eta}.
  \label{eq:raw-edge-increment-local}
\end{equation}
Using the preceding supremum estimate at the two points gives instead
\begin{equation}
  \int
  \abs{\mathcal K_rH(x)-\mathcal K_rH(y)}
  \leq
  C r^{-1+\eta/2}\norm H_{\mathcal H^\eta}.
  \label{eq:raw-edge-increment-global}
\end{equation}
If $d\leq\sqrt r$, divide~\eqref{eq:raw-edge-increment-local} by $d^\beta$; if $d>\sqrt r$, divide~\eqref{eq:raw-edge-increment-global} by $d^\beta$. Both cases give the factor $r^{-1+(\eta-\beta)/2}$. The supremum estimate is bounded by the same scale after changing the constant by a factor depending on $T$. Since $\eta-\beta>0$, the power of $r$ is integrable at zero.
\end{proof}

If $H_i:\Torus\times E_i\to\R$ and $m_i$ are finite positive measures, then the product field satisfies
\begin{equation}
  \norm{H_1H_2}_{\mathcal H^\eta(m_1\otimes m_2)}
  \leq
  \norm{H_1}_{\mathcal H^\eta(m_1)}
  \norm{H_2}_{\mathcal H^\eta(m_2)}.
  \label{eq:raw-product-holder}
\end{equation}
Indeed the $L^1$ norm factorizes, and the spatial increment is split as
\begin{equation*}
  H_1(x)H_2(x)-H_1(y)H_2(y)
  =
  [H_1(x)-H_1(y)]H_2(x)
  +H_1(y)[H_2(x)-H_2(y)].
\end{equation*}

\begin{theorem}[Canonical raw signed measure for a fixed finite tree]
\label{thm:canonical-raw-measure}
Fix $0<\alpha<1$, $T>0$, $\lambda\in\R$, and assume
\begin{equation*}
  \phi''\in C^\alpha(\Torus).
\end{equation*}
For every finite planar full binary tree $\tau$ and every $0<\beta<\alpha$,
\begin{equation}
  \sup_{0\leq t\leq T}
  \norm{\mathscr R_\tau^{t,\cdot}}_{\mathcal H^\beta(\nu_\tau)}
  <\infty.
  \label{eq:raw-fixed-tree-holder}
\end{equation}
In particular, for every target $(t,x)$,
\begin{equation*}
  \mathscr R_\tau^{t,x}\in L^1(\nu_\tau).
\end{equation*}
Hence
\begin{equation}
  \boxed{
  \mu_\tau^{t,x}(A)
  =
  \int_A\mathscr R_\tau^{t,x}(\omega)
  \,d\nu_\tau(\omega),
  \qquad A\in\mathcal B(\Omega_\tau),
  }
  \label{eq:canonical-raw-measure}
\end{equation}
defines a finite countably additive signed measure. Its total variation is
\begin{equation*}
  \norm{\mu_\tau^{t,x}}_{\TV}
  =
  \int_{\Omega_\tau}
  \abs{\mathscr R_\tau^{t,x}}d\nu_\tau
  <\infty.
\end{equation*}
Moreover,
\begin{equation}
  \boxed{
  \mu_\tau^{t,x}(\Omega_\tau)
  =
  F_\tau(t,x).
  }
  \label{eq:canonical-raw-mass}
\end{equation}
No smallness condition is required.
\end{theorem}

\begin{proof}
We prove~\eqref{eq:raw-fixed-tree-holder} by induction on the number of internal vertices. For the leaf,~\eqref{eq:raw-density-leaf} and translation invariance give a uniform $\mathcal H^\eta(\gamma)$ bound for every $0<\eta\leq\alpha$.

Let $\tau=[\tau_1,\tau_2]$, fix $0<\beta<\alpha$, and choose
\begin{equation*}
  \eta=\frac{\alpha+\beta}{2},
  \qquad
  \beta<\eta<\alpha.
\end{equation*}
By the induction hypothesis, the two child raw fields have finite uniform $\mathcal H^\eta$ norms. Their product has finite $\mathcal H^\eta$ norm by~\eqref{eq:raw-product-holder}. Equation~\eqref{eq:raw-density-internal} is $\lambda\mathcal K_{t-s}$ applied to this product. Lemma~\ref{lem:raw-one-edge-loss} therefore gives
\begin{align*}
  \sup_{t\leq T}
  \norm{\mathscr R_\tau^{t,\cdot}}_{\mathcal H^\beta(\nu_\tau)}
  &\leq
  \abs\lambda C_{\beta,\eta,T}
  \int_0^T
  r^{-1+(\eta-\beta)/2}\,dr\\
  &\quad\times
  \sup_{s\leq T}
  \norm{\mathscr R_{\tau_1}^{s,\cdot}}_{\mathcal H^\eta}
  \sup_{s\leq T}
  \norm{\mathscr R_{\tau_2}^{s,\cdot}}_{\mathcal H^\eta}
  <\infty.
\end{align*}
Thus the density in~\eqref{eq:canonical-raw-measure} belongs to $L^1(\nu_\tau)$. Since $\nu_\tau$ is finite positive,~\eqref{eq:canonical-raw-measure} is automatically a finite countably additive signed measure.

It remains to prove~\eqref{eq:canonical-raw-mass}. For the leaf,
\begin{equation*}
  \mu_\bullet^{t,x}(\Omega_\bullet)
  =
  \int g(x+\sqrt t\xi)d\gamma(\xi)
  =
  P_tg(x)
  =
  F_\bullet(t,x).
\end{equation*}
Assume the identity for the two children. Absolute integrability of the root density permits Fubini. Integrating first over the child spaces in~\eqref{eq:raw-density-internal} gives
\begin{align*}
  \mu_\tau^{t,x}(\Omega_\tau)
  ={}&
  \lambda\int_0^t\int_\R
  \frac{\He_2(z)}{t-s}
  \Bigl[
    F_{\tau_1}(s,x_z)F_{\tau_2}(s,x_z)
  \notag\\
  &\hspace{7em}-
    F_{\tau_1}(s,x)F_{\tau_2}(s,x)
  \Bigr]
  d\gamma(z)ds.
\end{align*}
The mass functions inherit spatial H\"older regularity from~\eqref{eq:raw-fixed-tree-holder}. The centered Gaussian Hessian identity~\eqref{eq:centered-hessian-mark} therefore identifies the inner integral with
\begin{equation*}
  \partial_x^2P_{t-s}
  [F_{\tau_1}(s)F_{\tau_2}(s)](x).
\end{equation*}
The resulting time integral is exactly the deterministic recursion~\eqref{eq:tree-profile-recursion}.
\end{proof}

The strict loss in Lemma~\ref{lem:raw-one-edge-loss} is not an additional data hypothesis. For any fixed finite tree one may choose finitely many intermediate exponents below the original $\alpha$. What deteriorates with depth is the size of the constants. In particular, the sharp Banach-scale lower bound of order $(cn/\Delta)^n$ means that a stepwise first-moment proof cannot keep these finite-depth constants geometric while $n$ grows; it does not say that a fixed depth has infinite first moment.

\paragraph{Decorated maximal-left-patch skeletons.}
The raw measure above is indexed by the original planar binary tree. Later sections also index the same measure by its decorated maximal-left-patch skeleton. The decoration must record more than the chain lengths.

For a non-leaf tree, follow the maximal left-child chain from the root and suppose it contains $m$ internal vertices. Let
\begin{equation*}
  \sigma_1,\ldots,\sigma_m
\end{equation*}
be the successive right subtrees in chain order, from the root vertex to the bottom vertex. Define recursively
\begin{equation*}
  \mathsf D(\bullet)=\bullet,
  \qquad
  \mathsf D(\tau)
  =
  \pa{
    m;
    \mathsf D(\sigma_1),\ldots,\mathsf D(\sigma_m)
  }.
  \label{eq:decorated-patch-encoding}
\end{equation*}
Thus the decoration records both each patch length and the ordered side-attachment slot of every descendant patch.

\begin{proposition}[Tree--decorated-patch bijection]
\label{prop:tree-patch-bijection}
The map $\mathsf D$ in~\eqref{eq:decorated-patch-encoding} is a bijection between finite planar full binary trees and decorated maximal-left-patch skeletons.
\end{proposition}

\begin{proof}
Construct the inverse recursively. Given
\begin{equation*}
  \pa{m;D_1,\ldots,D_m},
\end{equation*}
build a chain of $m$ internal vertices connected through left children. Attach $\mathsf D^{-1}(D_j)$ as the right subtree at the $j$-th vertex, and make the left child of the bottom vertex a leaf. Recurse inside each $D_j$. This reconstruction is unique and is visibly inverse to~\eqref{eq:decorated-patch-encoding}.
\end{proof}

A contracted patch graph decorated only by chain lengths would not be injective, because it may forget the vertex of a long patch at which a side subtree is attached. The ordered side slots are therefore part of the term \emph{decorated patch skeleton} throughout the rest of the paper. Consequently tree-indexing and decorated-skeleton-indexing are only two labels for the same $\Omega_\tau$, $\nu_\tau$, and $\mu_\tau^{t,x}$.

Finally, for a right-comb tree the signed measure used in~\cref{sec:raw-faithful} is the restriction of~\eqref{eq:canonical-raw-measure} to the specified duration cylinder. The Fourier leaf projectors in that proof are compatible with the common-seed centered coupling in~\eqref{eq:raw-density-internal}: periodic heat-kernel translation invariance makes each projector a bounded function of the terminal Brownian increment and remaining time, independent of the starting point. Thus the comb lower bound is a lower bound for the canonical measures constructed here, not for a different coupling.


\subsection{Coarsening the raw mark space}

Let $\cC_\tau:\Omega_\tau\to\overline\Omega_\tau$ be measurable. The coarsened signed measure is
\begin{equation}
  \overline\mu_\tau
  =
  (\cC_\tau)_{\#}\mu_\tau.
  \label{eq:pushforward-def}
\end{equation}
Its total mass is preserved, while total variation satisfies
\begin{equation}
  \abs{F_\tau(t,x)}
  \leq
  \norm{\overline\mu_\tau}_{\TV}
  \leq
  \norm{\mu_\tau}_{\TV}.
  \label{eq:tv-coarsening-sandwich}
\end{equation}
If $\cC_\tau$ is constant, equality holds in the lower bound. Thus complete interior averaging is total-variation optimal among deterministic coarsenings of a fixed skeleton measure.

The identity map gives the canonical raw marked measure. A Gaussian-bridge coarsening keeps, along a consecutive Hessian chain, the total displacement
\begin{equation*}
  Y=\sum_j\sqrt{r_j}Z_j
\end{equation*}
and forgets the orthogonal bridge coordinates. The conditional identity~\eqref{eq:bridge-score} describes the induced signed score for a bare chain.

\subsection{Parabolic H\"older space}

Fix $0<\alpha<1$ and $T>0$. We use
\begin{equation*}
  X_{\alpha,T}
  =
  C^{\alpha/2,\alpha}([0,T]\times\Torus)
\end{equation*}
with norm
\begin{equation}
  \norm f_X
  =
  \norm f_\infty
  +[f]_{t;\alpha/2}
  +[f]_{x;\alpha}.
  \label{eq:X-norm}
\end{equation}
Standard heat-semigroup estimates imply that there is a finite constant $C_{\cD}(\alpha,T)$ such that
\begin{equation}
  \norm{\cD(fg)}_X
  \leq
  C_{\cD}(\alpha,T)\norm f_X\norm g_X.
  \label{eq:D-bilinear-bound}
\end{equation}
The exact value is not important below; what matters is that the constant is finite for fixed $(\alpha,T)$ and closes the quadratic Duhamel map. Standard parabolic Schauder theory used later may be found in~\cite{Evans2010,Friedman1964,Lieberman1996}.

\subsection{Coding-tree notation}

For the Nguwi--Penent--Privault construction, $H(\mathcal T_{t,x,c})$ denotes the multiplicative functional of the coding tree started at time $t$, position $x$, and root code $c$~\cite{NguwiPenentPrivault2023}. We use only two structural facts from the construction. First, an allowed composite code $g^*$ may be written
\begin{equation*}
  g
  =
  a\,\partial_{z_0}^{\lambda_0}\cdots
  \partial_{z_n}^{\lambda_n}f,
  \qquad a\neq0.
\end{equation*}
Second, for every active direction $j$, the mechanism set of a composite code contains a three-child Hessian branch whose first child is the code associated with $-\frac12\partial_{z_j}^2g$ and whose two side children carry the spatial derivative $\partial_x^{j+1}$. Repeating this branch along the first child creates the distinguished genealogy used in~\cref{sec:coding-tree}.

All coding-tree moment statements state their lifetime and mechanism-positivity assumptions locally. In particular, the negative theorem never requires a uniform positive lower bound because the selected proposal factors cancel their reciprocal compensators.
